\documentclass[11pt]{amsart}
\usepackage[margin=1in]{geometry}
\usepackage{amsmath,amssymb,amsthm}
\usepackage[hidelinks]{hyperref}

\newtheorem{theorem}{Theorem}

\DeclareMathOperator{\Iso}{Iso}
\DeclareMathOperator{\Age}{Age}

\title{Later Negative Answers to Two Guarded-Fraisse Questions in arXiv:2408.12755}
\author{FA Banach run packet}
\date{2026-06-23}

\begin{document}
\maketitle

\section*{Status}

This packet records a later-literature answer to two questions from
Cuth--de Rancourt--Doucha, \emph{Guarded Fraisse Banach spaces},
arXiv:2408.12755. In Section 7 the source asks:

\begin{quote}
Is every guarded Fraisse, respectively cofinally Fraisse, Banach space
\(\omega\)-categorical?
\end{quote}

\begin{quote}
Does every guarded Fraisse, respectively cofinally Fraisse, Banach space have
a closed age?
\end{quote}

The later paper of de Rancourt--Fakhoury, \emph{Isometric rigidity and
Fraisse properties of Orlicz sequence spaces}, arXiv:2604.02080, answers both
questions negatively, even in the cofinally Fraisse class.

\section*{Answering Theorem}

Let \(p>3\), and define the Orlicz function
\[
  M(t)=t^p e^{t-1}.
\]
Then the Orlicz sequence space \(\ell_M\) is cofinally Fraisse. Moreover,
\(\ell_M\) is not oligomorphic, does not contain an isomorphic copy of
\(\ell_2\), and has non-closed age.

\section*{Implications for arXiv:2408.12755}

The source paper uses \(\omega\)-categoricity in the continuous-logic sense,
equivalently oligomorphy of the isometry action on finite powers of the unit
ball. Therefore ``not oligomorphic'' is precisely a negative answer to the
\(\omega\)-categoricity question.

The answering paper proves that \(\ell_M\) is cofinally Fraisse, and cofinally
Fraisse spaces are guarded Fraisse. Since this same space is not oligomorphic,
Question \texttt{quest:GFImpliesOmegaCat} has a negative answer for both
``guarded'' and ``cofinally guarded'' versions.

The answering theorem also states that \(\ell_M\) has non-closed age. Since
\(\ell_M\) is cofinally Fraisse, Question \texttt{quest:closedAge} likewise
has a negative answer for both the guarded and cofinally Fraisse versions.

Finally, the source's broad request for new infinite-dimensional guarded or
cofinally Fraisse Banach spaces is answered positively by this Orlicz sequence
family. This does not settle the source's more specific subquestions about
\(L_p(L_q(L_r))\) or \(L_p(E)\).

\section*{Provenance}

The cheap run indexes contained no pre-existing promoted packet for this
answer when this note was created. The relevant source facts are:
\begin{itemize}
\item arXiv:2408.12755, Section 7, labels \texttt{quest:GFImpliesOmegaCat}
and \texttt{quest:closedAge}.
\item arXiv:2604.02080, abstract: the authors prove guarded Fraisse Orlicz
sequence spaces that are not \(\omega\)-categorical and have non-closed age,
and say this answers a question of Cuth--de Rancourt--Doucha.
\item arXiv:2604.02080, Corollary \texttt{cor:OrliczFraisse} and Theorem
\texttt{thm:Counterexample}: \(M(t)=t^p e^{t-1}\), \(p>3\), gives the stated
cofinally Fraisse counterexample.
\end{itemize}

\section*{References}

\begin{thebibliography}{9}

\bibitem{CDR}
M. Cuth, N. de Rancourt, and M. Doucha,
\emph{Guarded Fraisse Banach spaces}, arXiv:2408.12755.

\bibitem{RF}
N. de Rancourt and M. Fakhoury,
\emph{Isometric rigidity and Fraisse properties of Orlicz sequence spaces},
arXiv:2604.02080.

\end{thebibliography}

\end{document}
